\begin{lemma}[Lemma 5.1, with the decay rate of (4.17) of the thesis]
Let $n\ge2$, $p>1$ and let $\mu,\gamma_0,\gamma_1,\gamma_2,\gamma_3,\gamma_4>0$
with $\gamma_0\le\gamma_1$ (``the data'' of \S5), suppose the
Sobolev--Poincar\'e inequality on balls holds in dimension $n$
(\noderef{CitedSobolevPoincareBall}), and assume the cited Dirichlet result (C3)
with H\"older boundary data (\noderef{CitedDirichletHolder}) in the universal
form described below.  Then (A) and (B) below both hold.

\medskip
\noindent\textbf{(A) The estimate of Lemma 5.1.}  Let $M\ge0$, let $d>0$ and
let $\sigma\colon\R\to\R$.  Then there exist a nonnegative, bounded,
continuous, increasing function $\sigma_2$ with $\sigma_2(0)=0$ and a number
$R_1>0$, both depending only on the data, $M$, $d$ and $\sigma$, such that the
following holds.  Let $0\le\varepsilon\le1$, let $\Omega\subseteq\R^n$ be open,
let $\Omega'$ be open and bounded with $\overline{\Omega'}\subseteq\Omega$ and
$d\le d(\Omega',\partial\Omega)$ (i.e.\ $d\le|x-y|$ for every $x\in\Omega'$ and
every $y\notin\Omega$).  Let $A$ and $B$ satisfy
(Q1) of the thesis--(Q6) of the thesis on $\Omega$ with these parameters and modulus
$\sigma$ (\noderef{EllipticStructureQ12}, \noderef{EllipticStructureQ345},
\noderef{EllipticStructureQ6}), and assume the cited H\"older estimates (C4)
(\noderef{CitedHolderOscillationSolution}).  Let
$u\in W^{1,p}(\Omega)\cap L^\infty(\Omega)$ with weak gradient $g=\nabla u$
(\noderef{MemW1pOn}) and $|u|\le M$ a.e.\ be a weak solution of (1.1) of the thesis
in $\Omega$ (\noderef{IsWeakSolutionOn}).  Let $x_0\in\Omega'$ and
$0<R\le R_1$ with $\overline B_R(x_0)\subseteq\Omega$, and let $v$ be the
comparison map for $u$ on $B_R(x_0)$ with the frozen operator
$h\mapsto A(x_0,U(x_0),h)$ (\noderef{FrozenComparisonOn}), i.e.\ the solution of
(5.2) of the thesis, with weak gradient $g_v$, continuous extension $V$ and
continuous representative $U$ of $u$ on $\overline B_R(x_0)$.  Put $w=u-v$, so
$\nabla w=g-g_v$.  Then
\[
\int_{B_R(x_0)}|\nabla w|^p\,dx\;\le\;\sigma_2(R)\int_{B_R(x_0)}(1+|\nabla u|^2)^{\frac p2}dx .
\]

\medskip
\noindent\textbf{(B) The decay rate.}  Let $0<\alpha\le1$ and $\delta>0$.  Then
there is $\delta'>0$, depending only on the data, $\alpha$ and $\delta$, with
the following property.  Let $M\ge0$, let $d>0$, let $\sigma\colon\R\to\R$, let
$c_\sigma>0$ and $C_1\ge0$, and suppose the modulus obeys the H\"older
condition that Theorem~Theorem~6 of the thesis assumes in addition (paper l.~1566--1567),
\[
\sigma(t)\le c_\sigma t^{\delta}\qquad(t\ge0).
\]
Then there is $R_1'>0$, depending only on the data,
$\alpha,\delta,c_\sigma,C_1,M,d$ and $\sigma$, such that the following holds.
Let $0\le\varepsilon\le1$, let $\Omega\subseteq\R^n$ be open, let $\Omega'$ be
open and bounded with $\overline{\Omega'}\subseteq\Omega$ and $d\le|x-y|$ for
every $x\in\Omega'$ and every $y\notin\Omega$.  Let $A$ and $B$ satisfy
(Q1) of the thesis--(Q6) of the thesis on $\Omega$ with these parameters and modulus
$\sigma$ (\noderef{EllipticStructureQ12}, \noderef{EllipticStructureQ345},
\noderef{EllipticStructureQ6}), and assume the cited H\"older estimates (C4)
(\noderef{CitedHolderOscillationSolution}).  Let
$u\in W^{1,p}(\Omega)\cap L^\infty(\Omega)$ with weak gradient $g=\nabla u$
(\noderef{MemW1pOn}) and $|u|\le M$ a.e.\ be a weak solution of (1.1) of the thesis
in $\Omega$ (\noderef{IsWeakSolutionOn}).  Let $x_0\in\Omega'$ and
$0<R\le R_1'$ with $\overline B_R(x_0)\subseteq\Omega$, and let $v,V,g_v,U$ be
as in (A), i.e.\ let the frozen comparison bundle
$\mathrm{FrozenComparisonOn}(n,p,x_0,R,A(x_0,U(x_0),\cdot),u,U,g,v,V,g_v)$ hold.
Suppose in addition that the representative $U$ is H\"older continuous of
exponent $\alpha$ on the closed ball with constant $C_1$:
\[
|U(x)-U(y)|\le C_1|x-y|^{\alpha}
\qquad\bigl(x,y\in\overline B_R(x_0)\bigr).
\]
Then
\[
\int_{B_R(x_0)}|\nabla w|^p\,dx\;\le\;R^{\delta'}\int_{B_R(x_0)}(1+|\nabla u|^2)^{\frac p2}dx .
\]

\medskip
\noindent\textbf{Why (B) is a second clause and not a strengthening of (A), and
why $\alpha$ and $\delta$ are bound inside it.}  Lemma 5.1 as the thesis states
it at l.~1531--1537 assumes no rate for the modulus of (Q6) of the thesis: ``there
exist a bounded continuous increasing function $\sigma_2(t)$ such that
$\sigma_2(0)=0$ and $R_1>0$ depending only on the data and
$d(\Omega',\partial\Omega)$''.  Neither $\alpha$ nor $\delta$ occurs there, and
$\delta$ is the exponent of the H\"older condition that only
Theorem~Theorem~6 of the thesis adds.  So $\alpha$ and $\delta$ must not be in scope in
(A): they are bound at the head of (B) only, and $\sigma_2$ and $R_1$
accordingly depend on exactly the quantities the thesis lists ($\sigma$ being
named because the modulus is a parameter of this lemma rather than fixed once
and for all).  Theorem 5 (\noderef{SolutionMorreyRegularity}) consumes (A) and
nothing else; Theorem 6 consumes (B).

\medskip
\noindent\textbf{On the dependence of $\delta'$.}  Theorem~Theorem~6 of the thesis says
$\lambda_0$ and $\beta=\lambda_0/p$ ``depend exclusively on the data'', while
the \emph{norm} depends also on $d(\Omega',\partial\Omega)$ and
$\|u\|_{W^{1,p}(\Omega)}$; its \S4 counterpart states the same list explicitly
at l.~1359--1361 and adds $\delta$ (``for a certain $\delta'>0$, depending only
on $\delta,\lambda,\Lambda,\nu_0,\nu_1,n,p$ and $a$'').  So $\delta'$ depends on
the data, on $\delta$ and on $\alpha$ --- $\alpha$ being ours rather than the
thesis', because we make the H\"older exponent a bound parameter --- and must
not depend on $M$, on $d$, on $\sigma$ or on $c_\sigma$; it is accordingly bound
before all four.  The radius $R_1'$ is bound after $c_\sigma$ and $C_1$ and may
depend on everything.

\medskip
\noindent\textbf{Why the H\"older exponent $\alpha$ is a parameter and the
H\"older datum a hypothesis.}  The thesis derives $\sigma_0$ and $\sigma_1$ ---
printed at l.~1524--1528 --- from a bound
$\sup_{B_R}|u(x_0)-u(x)|\le C_1R^{\alpha}$ in which $\alpha$ is the exponent of
the cited H\"older estimates (C4).  That exponent is \emph{not} available here
with the dependence $\delta'$ needs: (C4) (\noderef{CitedHolderOscillationSolution})
binds $M$ before its exponent, so an $\alpha$ extracted from (C4) inside this
lemma would be a function of $M$, and so would $\delta'$.  The paper's own route
avoids this, because Theorem~Theorem~6 of the thesis is stated in the situation of
Theorem~Theorem~5 of the thesis and Theorem 5's part (iii)
(\noderef{SolutionMorreyRegularity}) gives $u\in C^\eta$ for \emph{every}
$\eta<1$: a consumer may therefore fix $\alpha$ --- $\alpha=\tfrac12$ serves ---
and absorb the resulting constant into $C_1$.  Accordingly $\alpha$ is bound at
the head of (B), the $C^\alpha$ datum appears only inside (B), and $C_1$ is
bound with $c_\sigma$, before $R_1'$ and after $\delta'$.  The datum is imposed
on the continuous representative $U$ and not on $u$: $u$ is a $W^{1,p}$
function, so a pointwise H\"older bound on it is not what Theorem 5 (iii)
delivers, whereas \noderef{FrozenComparisonOn} already designates $U$ as the
continuous representative and this lemma already uses the value $U(x_0)$ in the
frozen operator.

\medskip
\noindent\textbf{Two cited results, and why they are at the head.}  The proof
instantiates Lemma~Lemma~3.1 of the thesis (\noderef{DirichletProblemRestricted}) at the
normalised frozen operator, and that node requires two cited results of the
paper's Note: the Sobolev--Poincar\'e inequality on balls (C6)
(\noderef{CitedSobolevPoincareBall}), which is a statement about $n$ alone, and
the cited Dirichlet result (C3) with H\"older boundary data
(\noderef{CitedDirichletHolder}), which is carried here for \emph{every} field
$A'\colon\R^n\to\R^n$ and all parameters $\varepsilon',\gamma_0',\gamma_1'$ with
$\gamma_0'>0$ and $\gamma_1'>0$.  The universal shape is forced: the field (C3)
is applied to is the normalised frozen operator $\widetilde A=A^0-A^0(0)$, which
is produced inside the proof at a centre and a value chosen there, so no single
instance can be named in the statement.  Both clauses are propositions
mentioning only $n$ and $p$, both of which are already bound above them, and
each is universally quantified over its remaining parameters; so placing them at
the head leaves the dependence lists of $\sigma_2$, $R_1$, $\delta'$ and $R_1'$
exactly as stated above, and no ``depending only on'' phrase changes.  This is
the opposite of what would happen if a genuine \emph{parameter} were hoisted:
$\alpha$ and $\delta$ are real numbers that the constants may be functions of,
which is why they are bound inside clause (B) and not at the head.
\end{lemma}
\begin{proof}
Throughout, $B_R$ abbreviates $B_R(x_0)$, $A^0$ the frozen field
$h\mapsto A(x_0,U(x_0),h)$, $\widetilde A$ its normalisation
$h\mapsto A^0(h)-A^0(0)$, and $w$ the difference $u-v$, whose weak gradient on
$B_R$ is $g-g_v$.  Write
\[
W=\int_{B_R}\bigl(\varepsilon+|g|^2+|g_v|^2\bigr)^{\frac p2-1}|g-g_v|^2dx,
\qquad
J=\int_{B_R}\bigl(1+|g|^2\bigr)^{\frac p2}dx,
\qquad
N=\int_{B_R}|g-g_v|^p\,dx .
\]

\medskip
\noindent\emph{Step 0: the two constants fixed before everything else.}  Apply
\noderef{SolutionFrozenMonotonicityIntegral} with the present $n$, $p$ and
$r=\gamma_1/\gamma_0>0$ --- its hypotheses before the constant are $2\le n$,
$1<p$ and $0<r$ --- and let $C_M>0$ be the constant it produces; it depends only
on $n,p$ and $\gamma_1/\gamma_0$.  Apply \noderef{WeightedRemainderAbsorption}
with the parameter $p$, whose only hypothesis before the constants is $1<p$, and
let $b_0>0$ and $c>0$ be the threshold and the constant it produces, both
depending only on $p$.  Put $K=1/(\gamma_0C_M)>0$.

\medskip
\noindent\emph{Step 1: a split that keeps the quantifier order of both clauses
honest.}  Clause (A) asserts $\sigma_2$ and $R_1$, and clause (B) asserts
$R_1'$, \emph{before} the inner block in which
$\mathrm{EllipticStructureQ6}\;n\;p\;\Omega\;A\;\sigma$
(\noderef{EllipticStructureQ6}) and
$\mathrm{CitedHolderOscillationSolution}\;n\;p$
(\noderef{CitedHolderOscillationSolution}) are assumed; (C6) and (C3) were
hoisted to the head of the lemma, but those two were not.  Yet Step 8 below
builds all five properties of $\sigma_2$ out of the four $\sigma$-clauses of
\noderef{EllipticStructureQ6}, Step 7 builds $\alpha$ and $C_1$ out of
\noderef{CitedHolderOscillationSolution}, $R_1$ comes from continuity of
$\sigma_1$ at $0$ --- again a $\sigma$-clause consequence --- and (9c)'s $R_1'$
inherits that dependence.  The two are reconciled by splitting on propositions
that mention only quantities already bound.

Two such propositions are available.  First, the conjunction of the four
$\sigma$-clauses of \noderef{EllipticStructureQ6} --- \texttt{sigmabdd},
\texttt{sigmacont}, \texttt{sigmamono} and \texttt{sigmazero} --- which assert
respectively that $\sigma$ is bounded above, continuous, nondecreasing and zero
at $0$.  Each of the four mentions \emph{only} $\sigma$, and in both clauses
$\sigma$ is bound before the witnesses: in clause (A) before $\sigma_2$ and
$R_1$, in clause (B) before $R_1'$.  (The fifth clause, \texttt{Q6}, does
mention $\Omega$ and $A$ and is therefore \emph{not} decidable here; it is used
only inside the inner block, in Step 5's freezing estimate, where $\Omega$ and
$A$ are in scope.)  Second,
$\mathrm{CitedHolderOscillationSolution}\;n\;p$, a proposition in $n$ and $p$
alone --- both bound at the head of the lemma --- since all of its remaining
parameters are universally quantified inside it.

\emph{Clause \textup{(A)}.}  Decide the conjunction of the four
$\sigma$-clauses and of $\mathrm{CitedHolderOscillationSolution}\;n\;p$ before
choosing $\sigma_2$ and $R_1$.  If it \emph{fails}, then for every $\Omega$ and
$A$ the hypothesis
$\mathrm{EllipticStructureQ6}\;n\;p\;\Omega\;A\;\sigma$ is false, or
$\mathrm{CitedHolderOscillationSolution}\;n\;p$ is; either way the inner
implication of clause (A) has a false hypothesis and is satisfied by any choice,
and we take $\sigma_2\equiv0$ and $R_1=1$.  These meet the five demands the
statement makes: $R_1=1>0$; $\sigma_2$ is bounded (by $0$), continuous,
nondecreasing and nonnegative on all of $\R$; and $\sigma_2(0)=0$.  No
dependence list is affected, $0$ and $1$ depending on nothing at all.  If it
\emph{holds}, we fix the four $\sigma$-clauses and
$\mathrm{CitedHolderOscillationSolution}\;n\;p$ once here and run Steps 2--9
inside this branch; every constant produced below is then produced under them,
and their dependence lists are exactly as recorded, these hypotheses
contributing no new quantity.

\emph{Clause \textup{(B)}.}  Its rate exponent $\delta'$ needs no split: it is
produced by \noderef{ModulusPowerDecayTransfer} from $p,\alpha,\delta$ alone,
all bound before it.  Its radius $R_1'$ does need one, and the four
$\sigma$-clauses suffice: decide them before choosing $R_1'$, take $R_1'=1$ in
the negative branch --- where $\mathrm{EllipticStructureQ6}$ again fails for
every $\Omega$ and $A$, so the inner implication is vacuous --- and run Steps
2--9 in the positive branch.  Clause (B) does not need the second split: it
assumes its own H\"older datum at its own $\alpha$ and derives nothing from
$\mathrm{CitedHolderOscillationSolution}\;n\;p$, which it carries only because
the inner hypothesis list is shared with clause (A).

\medskip
\noindent\emph{Step 2: the two moduli, fixed before the data.}  Let
$\sigma$ be given with the four $\sigma$-clauses, say $\sigma\le C_\sigma$, and
let $\alpha$ and $C_1$ with $0<\alpha\le1$ and $C_1\ge0$ be given --- in clause
(A) they come from Step 7 below, in clause (B) they are hypotheses of the clause.
Apply \noderef{SolutionComparisonModulus} with the present $p$, that $\alpha$
and $C_1$, with $K=1/(\gamma_0C_M)$, with $\gamma_4$ and with $C_\sigma$ and
$\sigma$.  Its hypotheses are $1<p$, $0<\alpha$, $\alpha\le1$, $0\le C_1$,
$0<K$ (as $\gamma_0,C_M>0$), $0<\gamma_4$ and the four $\sigma$-clauses, all
present.  It produces $\sigma_0,\sigma_1$ that are nonnegative, bounded,
continuous and nondecreasing on all of $\R$ with
$\sigma_0(0)=\sigma_1(0)=0$, and that satisfy
\[
\sigma_0(t)=K\Bigl(\frac{p-1}p\,\sigma\bigl(t+C_1t^{\alpha}\bigr)
+4\gamma_4C_1t^{\alpha}\Bigr),
\qquad
\sigma_1(t)=K\cdot\frac1p\,\sigma\bigl(t+C_1t^{\alpha}\bigr)
\tag{9.1}
\]
for $0<t\le1$.  \emph{These are the thesis' own explicit formulas at
l.~1524--1528}, with its $C(p)$ written $C_M$ and with the constant $4$ in place
of its printed $2$ in the $\gamma_4$ term; the latter is Step 6's correction,
authorised by the source's adjudicated \texttt{CHECK} at l.~1514 and discussed
there.

Now fix the data $\varepsilon,\Omega,\Omega',A,B,u,U,g,x_0,R,v,V,g_v$ with all
the hypotheses of the inner block, together with $0<R\le1$ and the H\"older
datum
\[
|U(x)-U(x_0)|\le C_1R^{\alpha}\qquad\bigl(x\in\overline B_R\bigr).
\tag{9.2}
\]
From $\overline B_R\subseteq\Omega$ we get $B_R\subseteq\Omega$ and
$x_0\in\Omega$, and $|B_R|<\infty$ with $|B_R|>0$ (Preamble: the volume of a
ball).  Restricting the two $L^p$ memberships of \noderef{MemW1pOn} from
$\Omega$ to $B_R$ gives $u,g\in L^p(B_R)$, clause (2) of
\noderef{FrozenComparisonOn} gives $g_v\in L^p(B_R)$, and $g-g_v\in L^p(B_R)$ as
their difference.

\medskip
\noindent\emph{Step 3: the frozen field, before and after normalisation.}
Apply \noderef{FrozenSolutionStructureQ12} with the present
$n,p,\varepsilon,\gamma_0,\gamma_1$ --- its hypotheses are $2\le n$, $1<p$,
$0\le\varepsilon$, $0<\gamma_0$ and $0<\gamma_1$, the last from
$0<\gamma_0\le\gamma_1$ --- the open set $\Omega$, the field $A$, the modulus
$\sigma$, the point $x_1=x_0\in\Omega$ and the value $z_1=U(x_0)$, together with
the \noderef{EllipticStructureQ12} and \noderef{EllipticStructureQ6} hypotheses
of the present lemma.  It gives \emph{both}
\[
\mathrm{EllipticStructureQ12}\;n\;p\;\varepsilon\;\gamma_0\;\gamma_1\;
\mathrm{univ}\;A^0
\qquad\text{and}\qquad
\mathrm{EllipticStructureQ12}\;n\;p\;\varepsilon\;\gamma_0\;\gamma_1\;
\mathrm{univ}\;\widetilde A
\]
(\noderef{EllipticStructureQ12}).  \emph{Both are used, and for different
things}: the un-normalised bundle feeds Step 4's monotonicity, where only the
difference $A^0(a)-A^0(b)$ appears and the constant cancels, while the
normalised one feeds the dual memberships of Step 3 and the comparison of Step
7, each of which is stated for the normalised field.  This step is also where
pm-0034's gap is closed: the (Q1) and (Q2) clauses are almost-everywhere in $x$
and so say nothing at the single frozen point $x_0$ until upgraded, and the
upgrade --- which consumes (Q6) of the thesis --- is inside that node.

\medskip
\noindent\emph{Step 4: the Lebesgue data of the two pairings.}  Apply
\noderef{SolutionWeakFormDualData} with the present
$n,p,\varepsilon,\mu,\gamma_0,\dots,\gamma_4$, the open set $\Omega$, the field
$A$, the right-hand side $B$, the pair $(u,g)$ and the measurable set $B_R$.
Its hypotheses are $2\le n$, $1<p$, $0\le\varepsilon\le1$, $0<\gamma_0$,
$0<\gamma_1$, $0<\gamma_3$, $0<\gamma_4$, the two structure bundles
\noderef{EllipticStructureQ12} and \noderef{EllipticStructureQ345}, the
measurability of $B_R$, the inclusion $B_R\subseteq\Omega$, $|B_R|<\infty$ and
$u,g\in L^p(B_R)$, all present or recorded in Step 2.  It gives
\[
A(\cdot,u,g)\in L^{p/(p-1)}(B_R)
\qquad\text{and}\qquad
B(\cdot,u,g)\in L^1(B_R).
\]
Apply \noderef{FrozenFieldDualMemLp} twice, with the present
$n,p,\varepsilon,\gamma_0,\gamma_1$, with $A^0$ in the role of its $A$, with
$B_R$ in the role of its $S$, and with $g$ and then $g_v$ in the role of its
$g_v$.  Its hypotheses are $2\le n$, $1<p$, $0\le\varepsilon\le1$, $0<\gamma_0$,
$0<\gamma_1$, the bundle for the \emph{normalised} field $\widetilde A$ on
$\mathrm{univ}$ --- which is exactly the second conclusion of Step 3, so no
transfer is needed --- $|B_R|<\infty$, and the relevant $L^p(B_R)$ membership.
It gives
\[
A^0(g)\in L^{p/(p-1)}(B_R)
\qquad\text{and}\qquad
A^0(g_v)\in L^{p/(p-1)}(B_R).
\]
Each pairing of one of these three with $g-g_v\in L^p(B_R)$ is then integrable
on $B_R$ by H\"older's inequality at the conjugate exponents $\frac p{p-1}$ and
$p$ (Preamble), and that is what licenses every splitting of an integral below.

\medskip
\noindent\emph{Step 5: the comparison bound for $v$, and the essential bound on
$w$.}  Put $M_\pm=U(x_0)\pm C_1R^{\alpha}$.  On the sphere $S_R$ we have
$V=U$ by clause (8) of \noderef{FrozenComparisonOn}, and $S_R\subseteq\overline
B_R$, so (9.2) gives $M_-\le V\le M_+$ there.  Apply
\noderef{BallSolutionConstantBounds} with the present
$n,p,\varepsilon,\gamma_0,\gamma_1$, with $A^0$ in the role of its $A$, with the
centre $x_0$ and radius $R$, with $v$, $V$, $g_v$ in the roles of its $v$, $V$,
$g_v$, and with $M_-$, $M_+$ in the roles of its $m$, $M$.  Its hypotheses are
$2\le n$, $1<p$, $0\le\varepsilon$, $0<\gamma_0$, $\gamma_0\le\gamma_1$;
\noderef{CitedSobolevPoincareBall}, a hypothesis of the present lemma; $0<R$;
the bundle for the normalised field from Step 3; the weak-solution property
$\mathrm{IsWeakSolutionOn}\;n\;p\;B_R\;A^0\;0\;v\;g_v$
(\noderef{IsWeakSolutionOn}), which is clause (1) of
\noderef{FrozenComparisonOn} verbatim --- \emph{for $A^0$, not for
$\widetilde A$}, which is why no transfer of the weak form is needed;
$v\in W^{1,p}(B_R)$ with weak gradient $g_v$ (\noderef{MemW1pOn}), clause (2);
$V$ continuous on $\overline B_R$, clause (6); $V=v$ on $B_R$, clause (7); and
the boundary trapping just shown.  \emph{That node performs the normalisation
internally}, taking the bundle for $\widetilde A$ and the weak form for $A^0$,
so pm-0048's transfer argument is not needed here and is not reproduced.  Its
conclusion is
\[
M_-\le v\le M_+\quad\text{a.e. on }B_R,
\qquad\text{that is}\qquad
|v-U(x_0)|\le C_1R^{\alpha}\quad\text{a.e. on }B_R .
\tag{9.3}
\]
Since $U=u$ at every point of $B_R$ by clause (5) of
\noderef{FrozenComparisonOn}, (9.2) also gives
$|u(x)-U(x_0)|\le C_1R^{\alpha}$ for every $x\in B_R$.  Adding this to (9.3)
through the single value $U(x_0)$,
\[
|w|=|u-v|\le2C_1R^{\alpha}\le4C_1R^{\alpha}\qquad\text{a.e. on }B_R ,
\tag{9.4}
\]
the last step because $C_1R^{\alpha}\ge0$.  \emph{The route through
\noderef{DirichletProblemRestricted} that this node's earlier sketch names is
circular}: that node's uniqueness clause demands
$\mathrm{ContDiffOn}\;\R\;1\;v$ on $B_R$ and
$\mathrm{MemHolderOn}\;n\;\beta\;(\overline B_R)\;V$ of the competitor, which
\noderef{FrozenComparisonOn} does not supply, and
\noderef{CitedDirichletHolder} states its regularity only of the solution
\emph{it} produces, so using it to certify $v$ presupposes that $v$ is that
solution.  We compare $v$ with constants instead, which is what the citation
above does.  In particular the thesis' (5.3) of the thesis,
$\operatorname*{osc}(v,B_R)\le\operatorname*{osc}(u,B_R)$, is \emph{not} needed
as an intermediate: (9.3) bounds $v$ directly.

\medskip
\noindent\emph{Step 6: \textup{(5.4)}, \textup{(5.5)}, \textup{(5.7)} and
\textup{(5.8)}.}  Each of the four displays of the source's derivation is now a
single citation.

\emph{(5.4) of the thesis.}  Apply \noderef{SolutionFrozenTestIdentity} with the
present $n,p$, the open set $\Omega$, the field $A$, the right-hand side $B$,
the pair $(u,U,g)$, the centre $x_0$, the radius $R$, the bound
$T=4C_1R^{\alpha}$ and the comparison datum $(v,V,g_v)$.  Its hypotheses are
$1\le n$ (from $2\le n$), $1<p$, $\Omega$ open, $0<R$, $B_R\subseteq\Omega$, the
weak-solution property \noderef{IsWeakSolutionOn}, the comparison datum
\noderef{FrozenComparisonOn}, the three $L^{p/(p-1)}$ memberships and the $L^1$
membership from Step 4, $g-g_v\in L^p(B_R)$ from Step 2, $T\ge0$, and the
essential bound (9.4).  It gives
\[
I_1\ \le\ I_2+I_3,
\qquad
I_1=\int_{B_R}\bigl(A^0(g)-A^0(g_v),g-g_v\bigr)dx,
\]
with $I_2=\bigl|\int_{B_R}(A(x,u,g)-A^0(g),g-g_v)dx\bigr|$ and
$I_3=\bigl|\int_{B_R}B(x,u,g)\,w\,dx\bigr|$.

\emph{(5.5) of the thesis.}  Apply the conclusion of
\noderef{SolutionFrozenMonotonicityIntegral}, whose constant $C_M$ was fixed in
Step 0, at the present $\varepsilon,\gamma_0,\gamma_1$, the field $A^0$, the
centre $x_0$, the radius $R$ and the pair $(g,g_v)$.  Its remaining hypotheses
are $0\le\varepsilon$, $0<\gamma_0$, $\gamma_1\le(\gamma_1/\gamma_0)\gamma_0$ ---
an equality, since $\gamma_0\ne0$ --- $0<R$, the un-normalised bundle from Step
3, the three $L^p(B_R)$ memberships from Step 2, and the integrability of
$x\mapsto(A^0(g)-A^0(g_v),g-g_v)$, which is the H\"older pairing recorded at the
end of Step 4.  It gives
\[
C_M\gamma_0\,W\ \le\ I_1 .
\]

\emph{(5.6) of the thesis--(5.7) of the thesis.}  Apply
\noderef{SolutionFreezingErrorBound} with the present $n,p$, the open set
$\Omega$, the field $A$, the modulus $\sigma$, the function $u$, the fields $g$
and $g_v$, the point $x_0$, the value $z_0=U(x_0)$, the radius $R$ and the pair
$(\alpha,C_1)$.  Its hypotheses are $1\le n$, $1<p$, the bundle
\noderef{EllipticStructureQ6}, $x_0\in\Omega$, $0<R$, $B_R\subseteq\Omega$,
$0<\alpha$, $0\le C_1$, the pointwise bound
$|u(x)-U(x_0)|\le C_1R^{\alpha}$ on $B_R$ from Step 5, $g,g-g_v\in L^p(B_R)$,
and the two memberships $A(\cdot,u,g),A^0(g)\in L^{p/(p-1)}(B_R)$ from Step 4.
It gives
\[
I_2\ \le\ \sigma\bigl(R+C_1R^{\alpha}\bigr)
\Bigl\{\frac{p-1}p\,J+\frac1p\,N\Bigr\}.
\]

\emph{(5.8) of the thesis.}  Apply \noderef{SolutionRightHandSideBound} with the
present $n,p,\varepsilon,\mu,\gamma_2,\gamma_3,\gamma_4$, the open set
$\Omega$, the field $A$, the right-hand side $B$, the pair $(u,v)$, the field
$g$, the centre $x_0$, the radius $R$ and the pair $(\alpha,C_1)$.  Its
hypotheses are $1\le n$, $1<p$, $0<\gamma_4$, $0<R$, $B_R\subseteq\Omega$,
$0<\alpha$, $0\le C_1$, the bundle \noderef{EllipticStructureQ345},
$g\in L^p(B_R)$, the integrability $B(\cdot,u,g)\in L^1(B_R)$ from Step 4, and
the essential bound (9.4).  It gives
\[
I_3\ \le\ 4C_1\gamma_4R^{\alpha}\,J .
\]
\emph{On the constant $4$}: the source's printed (5.8) of the thesis carries
$2C_1\gamma_4$, and its adjudicated \texttt{CHECK} at l.~1514 records that the
printed chain for $|w|$ --- through a point of $\partial B_R$ --- really yields
$4C_1R^{\alpha}$.  Step 5 above reaches the \emph{sharper} bound
$2C_1R^{\alpha}$, because (9.3) lets $u$ and $v$ be compared through the single
value $U(x_0)$ rather than through a boundary point; we weaken it to
$4C_1R^{\alpha}$ in (9.4) and carry $4$ into $\sigma_0$, which is the constant
the absorption step below and the moduli (9.1) are written with.  Nothing in the
proof asserts the printed inequality, and the slack is harmless because
$\sigma_0$ absorbs it.

\medskip
\noindent\emph{Step 7: \textup{(5.9)}.}  Chain the four displays of Step 6:
\[
C_M\gamma_0\,W\ \le\ I_1\ \le\ I_2+I_3\ \le\
\sigma\bigl(R+C_1R^{\alpha}\bigr)\Bigl\{\frac{p-1}p\,J+\frac1p\,N\Bigr\}
+4C_1\gamma_4R^{\alpha}\,J ,
\]
and regroup the right-hand side as
$\bigl(\frac{p-1}p\sigma(R+C_1R^{\alpha})+4\gamma_4C_1R^{\alpha}\bigr)J
+\bigl(\frac1p\sigma(R+C_1R^{\alpha})\bigr)N$.  Multiplying by
$K=1/(\gamma_0C_M)>0$ and reading off (9.1) at $t=R$, which is legitimate
because $0<R\le1$, gives the thesis' (5.9) of the thesis,
\[
W\ \le\ \sigma_0(R)\,J+\sigma_1(R)\,N .
\tag{9.5}
\]

\medskip
\noindent\emph{Step 8: the absorption, and clause \textup{(A)}.}  Since
$\sigma_1$ is continuous with $\sigma_1(0)=0$ and $b_0>0$, there is $\rho>0$
with $\sigma_1(\rho)\le b_0$.  In clause (A), take
\[
\sigma_2(t)=c\bigl(\sigma_0(t)+\sigma_0(t)^{\frac p2}+\sigma_1(t)^{\frac p2}\bigr),
\qquad
R_1=\min\{R_0,1,\rho\},
\]
where $R_0$ is the radius produced in Step 9 below.  The five demands the
statement makes of $\sigma_2$ follow from the corresponding properties of
$\sigma_0$ and $\sigma_1$ from Step 2, together with the facts that
$s\mapsto s^{p/2}$ is continuous, nondecreasing and nonnegative on $[0,\infty)$
with $0^{p/2}=0$ (pm-0024, $p/2>0$), and that $c>0$: boundedness, because
$\sigma_0,\sigma_1$ are bounded by some $C_b\ge0$ and then
$\sigma_2\le c(C_b+2C_b^{p/2})$; continuity and monotonicity, as a nonnegative
combination of continuous nondecreasing functions; $\sigma_2(0)=0$; and
$\sigma_2\ge0$.

Now let $0<R\le R_1$.  Then $R\le R_0$, $R\le1$ and $R\le\rho$, so
$\sigma_1(R)\le\sigma_1(\rho)\le b_0$ by monotonicity.  Apply the conclusion of
\noderef{WeightedRemainderAbsorption}, whose constants were fixed in Step 0,
with the dimension $n$, the present $\varepsilon$, with $\sigma_0(R)$ in the
role of its $a$ and $\sigma_1(R)$ in the role of its $b$, with $B_R$ in the role
of its $B$, and with $g$ and $g_v$ in the roles of its $G$ and $G_v$.  Its
hypotheses are $0\le\varepsilon\le1$; $0\le\sigma_0(R)$ and
$0\le\sigma_1(R)\le b_0$; $B_R$ measurable with $|B_R|\ne\infty$ (Preamble: the
volume of a ball); the three $L^p(B_R)$ memberships of $g$, $g_v$ and $g-g_v$
from Step 2; and its weighted hypothesis, which is exactly (9.5).  Its
conclusion is
\[
N\ \le\ c\bigl(\sigma_0(R)+\sigma_0(R)^{\frac p2}+\sigma_1(R)^{\frac p2}\bigr)J
=\sigma_2(R)\,J ,
\]
which is clause (A).  \emph{The single closed form the thesis displays at the
end of its two-regime argument is the $1<p<2$ branch only and is false for
$p>2$}; that is reported and repaired inside
\noderef{WeightedRemainderAbsorption}, which owns it, and is not reintroduced
here.  The paper's own proof of Lemma~Lemma~5.1 of the thesis is the one sentence at
l.~1540--1542, ``this follows from (5.9) of the thesis like Lemma~Lemma~4.1 of the thesis
followed from (4.7) of the thesis''; the adjudicated \texttt{CHECK} immediately above
it records that the analogue of (5.9) of the thesis is (4.8) of the thesis and not
(4.7) of the thesis, so the step cited here is the passage from (5.9) of the thesis
exactly as Lemma~Lemma~4.1 of the thesis passes from (4.8) of the thesis.

Finally $\sigma_2$ and $R_1$ depend only on
$n,p,\mu,\gamma_0,\dots,\gamma_4,M,d$ and $\sigma$, because $\sigma_0$,
$\sigma_1$, $C_1$, $\alpha$ and $R_0$ do and because $c$, $b_0$ and $C_M$ depend
only on $n,p$ and $\gamma_1/\gamma_0$.  There is \emph{no} $\alpha$ and
\emph{no} $\delta$ in that list as free parameters: $\alpha$ is produced from the
data in Step 9 and is not a parameter of clause (A), and $\delta$ appears only
in clause (B).

\medskip
\noindent\emph{Step 9: the two H\"older data, and clause \textup{(B)}.}

\textbf{(9a) Clause \textup{(A)}'s $\alpha$, $R_0$ and $C_1$.}  Apply
\noderef{SolutionHolderDatumOnBall} with the present
$n,p,\mu,\gamma_0,\dots,\gamma_4$, with $M$ and with $d$.  Its hypotheses are
$1\le n$, $1<p$, $0<\mu$, $0<\gamma_0$, $\gamma_0\le\gamma_1$, $0<\gamma_2$,
$0<\gamma_3$, $0<\gamma_4$, $0\le M$, $0<d$ and
$\mathrm{CitedHolderOscillationSolution}\;n\;p$, which Step 1's split has fixed
in this branch.  It produces $\alpha,R_0,C_1$ with $0<\alpha\le1$, $R_0>0$ and
$C_1\ge0$, depending only on $n,p,\mu,\gamma_0,\dots,\gamma_4,M$ and $d$, and
\emph{bound before $\varepsilon$ and before all the data}, which is what clause
(A) requires of them: $\sigma_2$ inherits $R^{\alpha}$ through (9.1), and clause
(A) binds $\sigma_2$ and $R_1$ before $\varepsilon$.  An $\alpha$ depending on
$\varepsilon$ would require one $\sigma_2$ to dominate $R^{\alpha(\varepsilon)}$
for every $\varepsilon$, and since $R^{\alpha}\to1$ as $\alpha\downarrow0$ for
each fixed $R\in(0,1)$, no nondecreasing function with $\sigma_2(0)=0$ could do
that.  Given the data and $0<R\le R_1\le R_0$, the node's remaining hypotheses
are $0\le\varepsilon\le1$, $\Omega$ open, the separation
$d\le\mathrm{dist}(x,y)$ for $x\in\Omega'$ and $y\notin\Omega$, $x_0\in\Omega'$,
the two structure bundles, the weak-solution property, $|u|\le M$ a.e.\ on
$\Omega$, $0<R$, $R\le R_0$, the continuity of $U$ on $\overline B_R$ (clause
(4) of \noderef{FrozenComparisonOn}) and $U=u$ on $B_R$ (clause (5)).  Its
conclusion is exactly (9.2), which is what Steps 2 and 5--6 consumed.

\textbf{(9b) Clause \textup{(B)}'s H\"older datum.}  In clause (B) the datum is
a hypothesis rather than a consequence: it reads
$|U(x)-U(y)|\le C_1\mathrm{dist}(x,y)^{\alpha}$ for $x,y\in\overline B_R$.  Read
at $y=x_0$, and using $\mathrm{dist}(x,x_0)\le R$ together with monotonicity of
$s\mapsto s^{\alpha}$ on $[0,\infty)$ and $C_1\ge0$, it gives (9.2) directly.
Nothing in Steps 2--8 uses any property of $\alpha$ and $C_1$ beyond
$0<\alpha\le1$, $C_1\ge0$ and (9.2) itself, all of which clause (B) supplies, so
the whole chain runs again on clause (B)'s own data.  \emph{What is combined with
it is the absorption inequality of Step 8, re-run on clause \textup{(B)}'s data
--- not the conclusion of clause \textup{(A)}, which is a different statement at
a different threshold and is not used here.}

\textbf{(9c) The decay rate.}  Let $0<\alpha\le1$ and $\delta>0$ be given.
Apply \noderef{ModulusPowerDecayTransfer} with the present $p$, that $\alpha$
and that $\delta$; its hypotheses are $1<p$, $0<\alpha$, $\alpha\le1$ and
$0<\delta$.  It produces $\delta'>0$ depending only on $p,\alpha,\delta$, which
is the binder order clause (B) requires.  Now let $M,d,\sigma,c_\sigma,C_1$ and
the rate hypothesis $\sigma(t)\le c_\sigma t^{\delta}$ for $t\ge0$ be given,
with $0\le M$, $0<d$, $0<c_\sigma$ and $0\le C_1$, and let $\sigma_0,\sigma_1$
be the moduli of Step 2 for this $\alpha$ and $C_1$.  Apply the node's
conclusion at $c_\sigma$, $C_1$, $K=1/(\gamma_0C_M)$, $\gamma_4$ and $c$, and at
$\sigma,\sigma_0,\sigma_1$: its remaining hypotheses are $0<c_\sigma$,
$0\le C_1$, $0<K$, $0<\gamma_4$, $0<c$, the rate hypothesis, the two closed
forms (9.1) on $(0,1]$, and the nonnegativity of $\sigma_0$ and $\sigma_1$, all
in hand.  It produces $R_2$ with $0<R_2\le1$ and
\[
\sigma_2(t)=c\bigl(\sigma_0(t)+\sigma_0(t)^{\frac p2}
+\sigma_1(t)^{\frac p2}\bigr)\ \le\ t^{\delta'}
\qquad\bigl(0<t\le R_2\bigr).
\]
Take $R_1'=\min\{R_2,\rho\}>0$, with $\rho$ as in Step 8.  For $0<R\le R_1'$ we
have $R\le R_2\le1$ and $R\le\rho$, so Step 8's absorption applies and gives
$N\le\sigma_2(R)J$, while $J\ge0$ because its integrand is nonnegative.  Hence
\[
N\ \le\ \sigma_2(R)\,J\ \le\ R^{\delta'}J ,
\]
which is clause (B).  The whole dependence on $c_\sigma,C_1,M,d$ and $\sigma$ has
been pushed into the threshold $R_1'$ --- the thesis' ``a radius threshold that
does depend on $c$'' --- while $\delta'$ depends only on $p,\alpha$ and
$\delta$.
\end{proof}
